\documentclass[11pt, a4paper]{article}

% ============================================================================
% PACKAGES
% ============================================================================
\usepackage[margin=1in]{geometry}
\usepackage{amsmath, amssymb}
\usepackage{booktabs}
\usepackage{graphicx}
\usepackage{hyperref}
\usepackage{xcolor}
\usepackage{float}
\usepackage{enumitem}
\usepackage{caption}
\usepackage{subcaption}
\usepackage{fancyhdr}
\usepackage{titlesec}
\usepackage{parskip}
\usepackage{natbib}
\bibliographystyle{plainnat}

% ============================================================================
% STYLING
% ============================================================================
\hypersetup{colorlinks=true, linkcolor=black, urlcolor=blue, citecolor=black}
\pagestyle{fancy}
\fancyhf{}
\rhead{Queen's Quant Club}
\lhead{Clustering-Driven Pair Discovery}
\rfoot{Page \thepage}
\renewcommand{\headrulewidth}{0.4pt}

\titleformat{\section}{\Large\bfseries}{}{0em}{}[\vspace{-0.5em}\rule{\textwidth}{0.4pt}]
\titleformat{\subsection}{\large\bfseries}{}{0em}{}

% ============================================================================
% TITLE
% ============================================================================
\title{
    \vspace{-1cm}
    \textbf{Clustering-Driven Pair Discovery:\\From Semiconductors to Cross-Sector Markets} \\[0.3cm]
    \large Using Unsupervised Learning to Identify and Validate \\
    Transient Mean-Reverting Relationships Across 142 Tickers
}
\author{
    Queen's Quant Club
}
\date{February 2026}

% ============================================================================
% DOCUMENT
% ============================================================================
\begin{document}
\maketitle
\thispagestyle{fancy}

% ============================================================================
% ABSTRACT
% ============================================================================
\begin{abstract}
We present a two-phase research project on transient correlation detection using unsupervised clustering for pairs trading. In Phase~1, three clustering algorithms---OPTICS, DBSCAN, and KMeans---are applied to 40 semiconductor stocks over 12 months of hourly data, converging on 8 consensus pairs that reflect genuine subsector structure (equipment, RF, analog, EDA). Of 26 tradeable pairs, 14 (54\%) produce positive out-of-sample P\&L under a z-score strategy.

In Phase~2, the framework is extended to 142 tickers across 5 sectors (Technology, Healthcare, Energy, Financial Services, Industrials) using a systematic three-layer screening pipeline. We replace the cointegration-based validation (0\% pass rate on transient correlations) with a 5-test scored framework (ADF, half-life, Hurst exponent, variance ratio, rolling correlation stability), yielding 3,643 scored pairs with 605 strong and 1,543 moderate. Enhanced backtesting with adaptive z-score optimization, Kalman hedge ratios, and 10bps transaction costs shows 64\% of pairs profitable under the Kalman strategy. Walk-forward validation confirms top pairs (AFRM--ZETA Sharpe 2.07, ENPH--RUN 1.82), and permutation testing identifies 1,148 of 3,539 pairs as statistically significant at $Z > 1.645$ (one-sided $p < 0.05$). Cross-sector discoveries include RTX--SHEL (Industrials/Energy) and a crypto mining cluster (APLD/HUT/CIFR/IREN spanning Technology and Financial Services).
\end{abstract}

\tableofcontents
\newpage

% ============================================================================
% 1. INTRODUCTION
% ============================================================================
\section{Introduction}

Pairs trading is among the oldest systematic strategies in quantitative finance \cite{gatev2006, vidyamurthy2004}, yet the first step---\textit{which} pairs to trade---is often approached naively. Traditional methods screen for cointegration \cite{engle1987} across all $\binom{N}{2}$ combinations, which scales poorly and produces many spurious results.

We propose an alternative: use unsupervised clustering on real-time feature data to let the algorithm \textit{discover} which stocks are behaving similarly at any given moment. When two stocks consistently land in the same cluster, it signals a relationship worth investigating with statistical tests.

This report documents a two-phase research program:

\begin{itemize}
    \item \textbf{Phase 1 (Semiconductors):} Proof of concept on 40 hand-picked semiconductor tickers. Three clustering algorithms are compared, consensus pairs identified, and results validated out of sample with a z-score strategy.
    \item \textbf{Phase 2 (Cross-Sector):} Extension to 142 systematically screened tickers across 5 GICS sectors. The validation framework is rebuilt from scratch (replacing cointegration with a 5-test scored system), enhanced backtesting with Kalman hedge ratios and transaction costs is added, and statistical significance is established via permutation testing.
\end{itemize}

The key research questions are:

\begin{enumerate}
    \item Do clustering algorithms reliably identify structurally related pairs?
    \item Does the approach generalize beyond a single sector?
    \item What validation framework is appropriate for transient (non-cointegrated) correlations?
    \item Do discovered pairs survive honest out-of-sample testing with transaction costs?
    \item Are the clustering results statistically significant vs.\ random feature assignment?
\end{enumerate}

% ============================================================================
% 2. DATA & FEATURE ENGINEERING
% ============================================================================
\section{Data and Feature Engineering}

\subsection{Phase 1: Semiconductor Universe}

Our initial universe consists of 40 semiconductor tickers spanning megacap leaders (NVDA, TSM, AVGO), equipment manufacturers (AMAT, LRCX, KLAC), specialized chip companies (QRVO, SWKS, NXPI), EDA vendors (CDNS, SNPS), and mid-caps (GFS, MRVL, CAMT). Hourly OHLCV data is sourced from Yahoo Finance, covering March 2025 to February 2026 ($\sim$1,580 hourly timestamps, $\sim$12 months).

The S\&P 500 (\texttt{\^{}GSPC}) is included as a market factor but excluded from clustering.

\subsection{Phase 2: Cross-Sector Universe}

To test generalizability, Phase~2 constructs a 142-ticker universe via a three-layer systematic screening pipeline:

\begin{enumerate}
    \item \textbf{Liquidity \& investability:} US-listed, market cap $> \$2$B, average daily volume $> 5$M shares, price $> \$5$.
    \item \textbf{Sector grouping:} Screens run per GICS sector---Technology, Healthcare, Energy, Financial Services, Industrials---yielding $\sim$40--80 stocks each, then pooled into one combined universe.
    \item \textbf{Fundamental quality:} Positive EBITDA (removes speculative pre-revenue companies whose price action is noise-driven).
\end{enumerate}

Data quality validation filters tickers with $<80\%$ non-null coverage or $>5\%$ zero-price observations.

\subsection{Feature Construction}

At each hourly timestamp, we compute 9 features per ticker:

\begin{table}[H]
\centering
\caption{Feature set for clustering.}
\begin{tabular}{@{}lll@{}}
\toprule
\textbf{Category} & \textbf{Feature} & \textbf{Window} \\
\midrule
Volatility      & Short-term rolling $\sigma$           & 50h ($\sim$1 week) \\
Market exposure & Rolling $\beta$ to S\&P 500           & 50h \\
Sector exposure & Rolling $\beta$ to leave-one-out sector avg & 50h \\
Momentum        & RSI                                    & 70h ($\sim$2 weeks) \\
Momentum        & 5-hour return                          & Instantaneous \\
Regime shift    & $\Delta\sigma$ (short $-$ medium vol)  & 50h vs.\ 147h \\
Regime shift    & $\Delta\beta_{\text{SPX}}$ (short $-$ medium) & 50h vs.\ 147h \\
Regime shift    & $\Delta\beta_{\text{sector}}$ (short $-$ medium) & 50h vs.\ 147h \\
Returns         & Hourly return                          & Instantaneous \\
\bottomrule
\end{tabular}
\label{tab:features}
\end{table}

The regime shift features (short-window minus medium-window) detect \textit{when} a stock's behaviour is changing relative to its recent baseline---the signal most relevant for transient relationship formation.

Before clustering, features are standardized (zero mean, unit variance) and reduced via PCA retaining 90\% of explained variance.

\subsection{Hedge Ratio Estimation}

Given price series $P_A$ and $P_B$, the OLS hedge ratio solves:
\[
\hat{\beta}_{\text{OLS}} = \frac{\text{Cov}(P_A, P_B)}{\text{Var}(P_B)}, \qquad \text{spread}_t = P_{A,t} - \hat{\beta} \cdot P_{B,t}
\]

The Kalman filter estimates a time-varying hedge ratio via the state-space model:
\begin{align}
\text{State:} \quad \boldsymbol{\theta}_t &= \boldsymbol{\theta}_{t-1} + \boldsymbol{w}_t, \quad \boldsymbol{w}_t \sim \mathcal{N}(\mathbf{0}, Q) \label{eq:kalman-state} \\
\text{Observation:} \quad P_{A,t} &= [P_{B,t}, \; 1] \cdot \boldsymbol{\theta}_t + v_t, \quad v_t \sim \mathcal{N}(0, R) \label{eq:kalman-obs}
\end{align}
where $\boldsymbol{\theta}_t = [\beta_t, \alpha_t]^\top$, $Q = 0.01 \cdot I_2$, and $R = 1.0$. In the terminal-beta approach, the filter runs on calibration data only; the final estimate $\hat{\beta}_T$ is applied as a fixed constant to out-of-sample data.

\subsection{Mean-Reversion Metrics}

\textbf{Half-life.} From the AR(1) regression $\Delta s_t = \phi \cdot s_{t-1} + c + \varepsilon_t$, the half-life of mean reversion is:
\[
h = -\frac{\ln 2}{\phi}, \qquad \phi < 0 \text{ (mean-reverting)}
\]

\textbf{Hurst exponent.} Using the variance ratio method, for lags $\tau \in \{2, \ldots, 20\}$:
\[
\sigma(\tau) = \text{std}(s_{t+\tau} - s_t), \qquad H = \text{slope of } \log \sigma(\tau) \text{ vs.\ } \log \tau
\]
$H < 0.5$ indicates mean reversion; $H = 0.5$ is a random walk; $H > 0.5$ is trending.

\subsection{Z-Score Trading Strategy}

The rolling z-score signal is:
\[
z_t = \frac{s_t - \bar{s}_{t,L}}{\hat{\sigma}_{t,L}}, \qquad \bar{s}_{t,L} = \frac{1}{L}\sum_{i=0}^{L-1} s_{t-i}
\]
where $L$ is the lookback window. Entry when $|z_t| \geq z_{\text{entry}}$; exit when $|z_t| \leq z_{\text{exit}}$. The enhanced optimizer performs grid search over $(z_{\text{entry}}, z_{\text{exit}}, L)$ on calibration data to maximize the calibration-period Sharpe ratio.

\subsection{Sharpe Ratio}

The annualized Sharpe ratio (used in Phase~2 enhanced backtesting) is computed from daily P\&L returns:
\[
\text{Sharpe} = \frac{\bar{r}_d}{\hat{\sigma}_d} \cdot \sqrt{252}
\]
where $\bar{r}_d$ and $\hat{\sigma}_d$ are the mean and standard deviation of daily returns from the cumulative P\&L curve. Pairs with fewer than 5 OOS trades receive $\text{Sharpe} = \text{NaN}$.

% ============================================================================
% 3. CLUSTERING METHODOLOGY
% ============================================================================
\section{Clustering Methodology}

\subsection{OPTICS}

OPTICS (Ordering Points To Identify the Clustering Structure) \cite{ankerst1999} is a density-based algorithm that does not require specifying the number of clusters \textit{a priori}. It identifies clusters of varying density and assigns a noise label ($-1$) to points that do not fit any cluster.

\textbf{Parameters:} \texttt{min\_samples=3}, \texttt{xi=0.05}, \texttt{min\_cluster\_size=3}.

OPTICS produces a high noise rate ($\sim$58\%), meaning most tickers at most timestamps are not assigned to any cluster. This is by design---it only forms clusters when there is genuine feature similarity.

\subsection{DBSCAN}

DBSCAN \cite{ester1996} uses a fixed density threshold ($\varepsilon$) with adaptive selection based on the k-nearest-neighbour distance distribution at each timestamp.

\textbf{Noise rate:} $\sim$28\%. Lower than OPTICS because the fixed threshold is less conservative.

\subsection{KMeans}

KMeans is centroid-based, assigning every point to a cluster. We select $k$ at each timestamp via silhouette score optimization and apply a post-hoc noise filter.

\textbf{Effective noise rate:} $\sim$10\% after filtering.

\subsection{Noise-Adjusted Co-Clustering Frequency}
\label{sec:noise}

For density-based algorithms, raw co-clustering frequency is misleading because many timestamps have one or both tickers in noise. We define:

\[
f_{\text{adj}}(A, B) = \frac{\text{co-cluster count}(A, B)}{|\{t : A \notin \text{noise at } t \text{ and } B \notin \text{noise at } t\}|}
\]

This measures: \textit{when both stocks are actually in real clusters, how often are they in the same one?}

\subsection{Algorithm Selection for Phase 2}

OPTICS was selected as the primary algorithm for Phase~2 based on Phase~1 findings:

\begin{itemize}
    \item Best at detecting variable-density clusters and transient relationship formations
    \item Conservative noise assignment means discovered pairs have high confidence
    \item Only 3 pairs appear in all 3 algorithms' top-20 (COP--DVN, CVX--XOM, SU--XOM, all energy)---consensus is weak, but OPTICS captures the broadest set of genuine relationships
\end{itemize}

% ============================================================================
% 4. PHASE 1: SEMICONDUCTORS
% ============================================================================
\section{Phase 1: Semiconductors}

\subsection{Algorithm Comparison}

\begin{table}[H]
\centering
\caption{Clustering quality metrics across algorithms (Phase 1).}
\begin{tabular}{@{}lccc@{}}
\toprule
\textbf{Metric} & \textbf{OPTICS} & \textbf{DBSCAN} & \textbf{KMeans} \\
\midrule
Valid clustering windows   & 1,362  & 1,579  & 1,579  \\
Avg clusters per timestamp & 3.3    & 1.5    & 3.1    \\
Avg noise rate             & 58\%   & 28\%   & $\sim$10\% \\
OOS stability ($r$)        & 0.671  & 0.832  & 0.627  \\
\bottomrule
\end{tabular}
\label{tab:algo-quality}
\end{table}

OOS stability is measured as the Pearson correlation of co-clustering frequencies between the first 67\% and last 33\% of timestamps.

\subsection{Consensus Pairs}

Despite fundamental differences, all three algorithms identify the same top pairs. Eight pairs appear in every algorithm's top-20 by co-clustering frequency:

\begin{table}[H]
\centering
\caption{Consensus pairs: appearing in all three algorithms' top-20.}
\begin{tabular}{@{}lcccc@{}}
\toprule
\textbf{Pair} & \textbf{OPTICS} & \textbf{KMeans} & \textbf{DBSCAN} & \textbf{Subsector} \\
\midrule
KLAC--LRCX  & 0.378 & 0.785 & 0.851 & Equipment \\
QRVO--SWKS  & 0.372 & 0.888 & 0.828 & RF \\
AMAT--LRCX  & 0.363 & 0.824 & 0.845 & Equipment \\
AMAT--KLAC  & 0.340 & 0.807 & 0.814 & Equipment \\
ADI--NXPI   & 0.279 & 0.816 & 0.900 & Analog \\
ADI--TXN    & 0.278 & 0.887 & 0.817 & Analog \\
NXPI--TXN   & 0.223 & 0.819 & 0.806 & Analog \\
ADI--SWKS   & 0.216 & 0.830 & 0.830 & Analog/RF \\
\bottomrule
\end{tabular}
\label{tab:consensus}
\end{table}

\subsection{Discovered Subsector Structure}

The algorithms---receiving only price-derived features, no fundamental data---independently discover the semiconductor industry's subsector structure:

\begin{itemize}
    \item \textbf{Equipment:} AMAT, LRCX, KLAC --- shared fab customers (TSMC, Samsung, Intel)
    \item \textbf{RF / 5G:} QRVO, SWKS, QCOM --- smartphone and 5G demand
    \item \textbf{Analog:} ADI, NXPI, TXN, MCHP --- industrial and automotive end markets
    \item \textbf{EDA:} CDNS, SNPS --- chip design software duopoly
\end{itemize}

This convergence between data-driven clustering and known industry structure provides strong evidence that discovered pairs reflect genuine economic relationships.

\subsection{Validation and Out-of-Sample Results}

Phase~1 used a 3-criteria binary validation: Engle-Granger cointegration \cite{engle1987} ($p < 0.05$), half-life \cite{chan2013} (5--60 days), and Hurst exponent \cite{hurst1951} ($< 0.5$). Of 33 filtered pairs, 4 passed all three and 22 near-missed (typically passing half-life and Hurst but failing cointegration).

All 26 tradeable pairs were backtested on the OOS window (last 33\%, $\sim$3.5 months) using a z-score mean-reversion strategy (lookback 20d, enter $|z| \geq 2.0$, exit $|z| \leq 0.5$, OLS hedge ratio):

\begin{table}[H]
\centering
\caption{Phase 1 top OOS performers.}
\small
\begin{tabular}{@{}lccccc@{}}
\toprule
\textbf{Pair} & \textbf{Coint $p$} & \textbf{Hurst} & \textbf{Trades} & \textbf{Sharpe} & \textbf{Win\%} \\
\midrule
NXPI--ON     & 0.271 & 0.413 & 2 & 3.55 & 100\% \\
CDNS--SNPS   & 0.620 & 0.447 & 3 & 2.08 & 100\% \\
QCOM--QRVO   & 0.019 & 0.205 & 2 & 1.44 & 100\% \\
KLAC--TSM    & 0.006 & 0.181 & 2 & 1.35 & 100\% \\
NXPI--POWI   & 0.568 & 0.372 & 5 & 0.85 & 80\% \\
\bottomrule
\end{tabular}
\label{tab:phase1-oos}
\end{table}

\textbf{14 of 26 tradeable pairs (54\%) produced positive OOS P\&L.} Trade counts are low (2--5), limiting statistical confidence.

\subsection{Look-Ahead Bias Lesson}

An earlier version estimated hedge ratios on the full data period and backtested on the same data, yielding 23/23 pairs profitable. After implementing the proper cal/OOS split, profitable pairs dropped to 14/26, and some strong in-sample pairs reversed (e.g., ADI--NXPI: IS Sharpe 1.08 $\to$ OOS Sharpe $-1.76$). This underscores the critical importance of out-of-sample validation.

\subsection{Consensus Does Not Predict OOS Performance}

\begin{table}[H]
\centering
\caption{Consensus vs.\ non-consensus OOS performance (Phase 1).}
\begin{tabular}{@{}lccc@{}}
\toprule
\textbf{Group} & \textbf{Tradeable} & \textbf{Avg Sharpe} & \textbf{Avg P\&L} \\
\midrule
Consensus (all 3 algorithms)     & 6  & $-0.251$ & $-13.73$ \\
Non-consensus                     & 20 & $+0.056$ & $-8.76$ \\
\bottomrule
\end{tabular}
\label{tab:consensus-oos}
\end{table}

Cross-algorithm consensus validates the \textit{relationship} but does not predict \textit{OOS profitability}.

% ============================================================================
% 5. PHASE 2: CROSS-SECTOR EXPANSION
% ============================================================================
\section{Phase 2: Cross-Sector Expansion}

\subsection{Universe and Clustering}

Phase~2 pools all 142 tickers into a single combined universe and clusters them together. Cross-sector pairs are discovered organically---the algorithm is not given sector labels. OPTICS identifies 7,836 co-clustering pairs across all timestamps. At the 0.08 noise-adjusted frequency threshold, 3,643 pairs are included for 5-test validation. Noise-adjusted frequency uses as denominator only the timestamps where \textit{both} tickers are non-noise (cluster\_id $\neq -1$).

\begin{table}[H]
\centering
\caption{Pair counts at the default noise-adjusted frequency threshold.}
\begin{tabular}{@{}lccc@{}}
\toprule
\textbf{Threshold} & \textbf{Total} & \textbf{Intra-Sector} & \textbf{Cross-Sector} \\
\midrule
\textbf{0.08 (default)} & \textbf{3,643} & \textbf{1,442} & \textbf{2,201} \\
\bottomrule
\end{tabular}
\label{tab:freq-thresholds}
\end{table}

The noise-adjusted frequency threshold of 0.08 determines which pairs enter the 5-test validation pipeline. At this threshold, 3,643 pairs qualify, with cross-sector pairs (2,201) outnumbering intra-sector (1,442).

\subsection{Why Cointegration Fails}

The Phase~1 validation used Engle-Granger cointegration as a primary criterion. Applied to the cross-sector universe, this produces a \textbf{0\% pass rate}. Cointegration tests for a long-run equilibrium, but transient correlations form and dissolve on shorter timeframes. Full-period stationarity tests fail by design.

Meanwhile, the Hurst exponent passes 90.5\% of pairs and half-life passes 82.3\%---the spreads \textit{do} mean-revert short-term. This mismatch between cointegration failure and mean-reversion success motivated the switch to a scored framework.

\subsection{5-Test Scored Validation Framework}

Each test scores 0 or 1; total score ranges 0--5.

\begin{table}[H]
\centering
\caption{5-test validation framework.}
\begin{tabular}{@{}llcl@{}}
\toprule
\textbf{Test} & \textbf{Threshold} & \textbf{Pass Rate} & \textbf{Rationale} \\
\midrule
ADF on spread \cite{dickey1979} & $p < 0.10$ & 38.3\% & Spread stationarity; 0.10 due to low ADF power \\
Half-life             & 5--60 days     & 82.3\% & Mean-reversion speed filter \\
Hurst exponent        & $H < 0.5$      & 90.5\% & Confirms mean-reverting behaviour \\
Variance ratio        & Reject RW @10\% & 21.8\% & Lo-MacKinlay \cite{lomackinlay1988}; hardest, detects short-run MR \\
Rolling corr stability & $> 0.5$       & 34.7\% & Returns correlation persistence across sub-windows \\
\bottomrule
\end{tabular}
\label{tab:5test}
\end{table}

\textbf{Classification:} Strong (4--5/5), Moderate (3/5), Weak (2/5), Fail ($<$2/5).

The real differentiators are ADF and variance ratio. A ``strong'' pair must pass at least one hard test plus most easy ones.

\subsection{Score Distribution}

\begin{table}[H]
\centering
\caption{Score distribution across 3,643 pairs.}
\begin{tabular}{@{}lcc@{}}
\toprule
\textbf{Classification} & \textbf{Count} & \textbf{Percentage} \\
\midrule
Strong (4--5)  & 605   & 16.6\% \\
Moderate (3)   & 1,543 & 42.4\% \\
Weak (2)       & 1,200 & 32.9\% \\
Fail (0--1)    & 295   & 8.1\%  \\
\midrule
\textbf{Tradeable ($\geq 3$)} & \textbf{2,148} & \textbf{59.0\%} \\
\bottomrule
\end{tabular}
\label{tab:score-dist}
\end{table}

\subsection{Intra-Sector vs Cross-Sector}

\begin{table}[H]
\centering
\caption{Intra-sector vs cross-sector comparison.}
\begin{tabular}{@{}lcccccc@{}}
\toprule
\textbf{Type} & \textbf{Pairs} & \textbf{Strong} & \textbf{Avg Score} & \textbf{Pass Rate} & \textbf{Avg Sharpe} & \textbf{Best Sharpe} \\
\midrule
Intra-sector   & 1,442 & 282 & 2.77 & 64.1\% & 1.95 & 1.96 \\
Cross-sector   & 2,201 & 323 & 2.62 & 55.6\% & 1.59 & 2.56 \\
\bottomrule
\end{tabular}
\label{tab:intra-cross}
\end{table}

Intra-sector pairs have higher average scores and pass rates, but cross-sector pairs produce the highest individual Sharpe (2.56).

\subsection{Sector Pair Matrix}

\begin{table}[H]
\centering
\caption{Top sector combinations by pair count.}
\small
\begin{tabular}{@{}llcccc@{}}
\toprule
\textbf{Sector 1} & \textbf{Sector 2} & \textbf{Pairs} & \textbf{Tradeable} & \textbf{Avg Score} & \textbf{Avg Sharpe} \\
\midrule
Technology    & Technology    & 880 & 504 & 2.64 & 1.93 \\
Energy        & Energy        & 270 & 230 & 3.09 & ---  \\
Industrials   & Technology    & 243 & 155 & 2.79 & 2.56 \\
Healthcare    & Technology    & 197 & 105 & 2.57 & 0.57 \\
Healthcare    & Healthcare    & 183 & 123 & 2.88 & ---  \\
Fin.\ Services & Technology  & 173 & 94  & 2.65 & ---  \\
\bottomrule
\end{tabular}
\label{tab:sector-matrix}
\end{table}

Energy--Energy pairs have the highest average score (3.09). Industrials--Technology produces the highest average Sharpe (2.56).

\subsection{Cross-Sector Discoveries}

\begin{table}[H]
\centering
\caption{Cross-sector pairs at 0.08 threshold.}
\begin{tabular}{@{}llcl@{}}
\toprule
\textbf{Pair} & \textbf{Sectors} & \textbf{Score} & \textbf{Notes} \\
\midrule
RTX--SHEL   & Industrials / Energy       & 4 (Strong)   & Defense--energy; genuinely surprising \\
APLD--HUT   & Technology / Fin.\ Services & 4 (Strong)   & Crypto mining infrastructure \\
APLD--CIFR  & Technology / Fin.\ Services & 3 (Moderate) & Crypto mining cluster \\
APLD--IREN  & Technology / Fin.\ Services & 3 (Moderate) & Crypto mining cluster \\
CPRT--CTSH  & Industrials / Technology   & 1 (Fail)     & Weak relationship \\
\bottomrule
\end{tabular}
\label{tab:cross-sector}
\end{table}

The APLD cluster (APLD/HUT/CIFR/IREN) consists of crypto mining infrastructure companies classified across Technology and Financial Services---a meaningful economic relationship. RTX--SHEL is the strongest cross-sector pair with a genuinely surprising defense--energy link.

% ============================================================================
% 6. ENHANCED BACKTESTING
% ============================================================================
\section{Enhanced Backtesting}

\subsection{Motivation}

The Phase~1 baseline strategy (static $z = 2.0$, OLS hedge, 20-day lookback) produces only 1--4 OOS trades per pair in $\sim$68 OOS days, yielding no valid Sharpe ratios. Phase~2 introduces adaptive parameters, Kalman hedge ratios, and transaction costs.

\subsection{Adaptive Z-Score Optimization}

Grid search on calibration data only:

\begin{itemize}
    \item Entry $z$: $\{1.0, 1.25, 1.5, 1.75, 2.0, 2.5\}$
    \item Exit $z$: $\{0.0, 0.25, 0.5, 0.75\}$
    \item Lookback: $\{10, 15, 20\}$ days
\end{itemize}

Scoring: maximize calibration Sharpe, requiring $\geq 5$ calibration trades. Returns \texttt{None} if no profitable parameter combination is found.

\subsection{Kalman Hedge Ratio}

The Kalman filter \cite{kalman1960} estimates the hedge ratio on calibration data, exponentially weighting recent observations for a superior terminal $\beta$ estimate.

\textbf{Critical fix:} An earlier version used an adaptive Kalman spread ($\beta$ updated at each OOS step), which creates artificial mean-reversion---the filter's prediction error has near-zero mean and no autocorrelation by construction, making z-score trading trivially profitable. The fixed-terminal-$\beta$ approach applies the calibration-estimated $\beta$ to OOS data, avoiding this pitfall.

\subsection{Three-Strategy Comparison}

\begin{table}[H]
\centering
\caption{Strategy comparison on top 50 tradeable pairs.}
\begin{tabular}{@{}llccc@{}}
\toprule
\textbf{Strategy} & \textbf{Description} & \textbf{Avg Trades} & \textbf{Profitable} & \textbf{Avg PnL} \\
\midrule
Baseline  & Static $z=2.0$, OLS, no costs & 1.9 & 54\% & 0.123 \\
Enhanced  & Optimized $z$, OLS, 10bps     & 2.8 & 57\% & 0.166 \\
Kalman    & Optimized $z$, Kalman $\beta$, 10bps & 2.6 & \textbf{64\%} & \textbf{0.895} \\
\bottomrule
\end{tabular}
\label{tab:strategy-comparison}
\end{table}

The Kalman strategy achieves the highest profitability (64\%) and average PnL (0.895), even with 10bps transaction costs. The enhanced optimizer generates more trades (2.8 vs 1.9 avg) by lowering entry thresholds.

\textbf{Why 50 of 2,148 tradeable pairs?} The enhanced backtest runs on the top 50 tradeable pairs by noise-adjusted frequency. The z-score optimizer requires $\geq 5$ calibration trades and a profitable calibration parameter set to produce optimized parameters. The Kalman strategy now runs independently of OLS optimization, falling back to default z-score params when the optimizer returns None.

\subsection{Top Enhanced Pairs}

\begin{table}[H]
\centering
\caption{Top pairs by enhanced PnL.}
\begin{tabular}{@{}llcccc@{}}
\toprule
\textbf{Pair} & \textbf{Type} & \textbf{Enh.\ PnL} & \textbf{Enh.\ Trades} & \textbf{Kalman PnL} \\
\midrule
AMAT--TSM   & intra & +30.30 & 2 & +54.19 \\
CVX--SU     & intra & +13.16 & 2 & +10.41 \\
MCHP--ON    & intra & +12.94 & 3 & +1.67  \\
COP--XOM    & intra & +11.23 & 6 & $-0.72$  \\
APLD--IREN  & cross & +10.40 & 2 & +9.47  \\
RTX--SHEL   & cross & +10.07 & 3 & +13.87 \\
HUT--RIOT   & intra & +9.91  & 8 & +5.65  \\
\bottomrule
\end{tabular}
\label{tab:top-enhanced}
\end{table}

\subsection{Walk-Forward Validation}

Walk-forward validation uses 5-split rolling calibration/OOS windows to test whether performance persists across multiple temporal splits:

\begin{table}[H]
\centering
\caption{Walk-forward results for top pairs.}
\begin{tabular}{@{}lcccc@{}}
\toprule
\textbf{Pair} & \textbf{Splits} & \textbf{Avg Sharpe} & \textbf{Std Sharpe} & \textbf{Total Trades} \\
\midrule
AFRM--ZETA  & 5 & \textbf{2.07} & 1.13 & 11 \\
ENPH--RUN   & 5 & \textbf{1.82} & 0.82 & 6  \\
COP--CVE    & 5 & 0.45          & 0.57 & 16 \\
\bottomrule
\end{tabular}
\label{tab:walk-forward}
\end{table}

AFRM--ZETA shows robust performance (Sharpe 2.07 with moderate variance) across 5 splits. ENPH--RUN has consistent Sharpe (1.82, std 0.82) reflecting strong solar energy sector coupling. COP--CVE has a lower Sharpe (0.45) but the highest trade count (16), providing more statistical confidence.

% ============================================================================
% 7. STATISTICAL SIGNIFICANCE
% ============================================================================
\section{Statistical Significance}

\subsection{Permutation Test Methodology}

We validate co-clustering significance using a feature-shuffle permutation test. At each of 80 sampled timestamps across 30 permutations, we shuffle feature vectors across tickers (preserving cross-feature correlation but breaking the ticker--feature mapping) and re-run the full StandardScaler $\to$ PCA $\to$ OPTICS pipeline. For each pair, we compute a $Z$-score comparing the observed co-clustering frequency to the null distribution.

\subsection{Results}

Of \textbf{3,539} pairs tested, \textbf{1,148 (32.4\%)} achieve $Z > 1.645$ (one-sided $p < 0.05$). Of these, \textbf{671} are also tradeable (score $\geq 3$). At the stricter $Z > 1.96$ threshold, 940 pairs (26.6\%) remain significant.

\begin{table}[H]
\centering
\caption{Top 10 statistically significant pairs by Z-score.}
\begin{tabular}{@{}lcc@{}}
\toprule
\textbf{Pair} & \textbf{Z-Score} & \textbf{Sector} \\
\midrule
PBR--PBR-A  & 38.59 & Energy \\
AAL--DAL    & 38.19 & Industrials \\
AAL--UAL    & 36.83 & Industrials \\
DAL--UAL    & 34.03 & Industrials \\
HAL--SLB    & 32.08 & Energy \\
KMI--WMB    & 30.65 & Energy \\
DVN--OXY    & 27.04 & Energy \\
CVX--XOM    & 26.86 & Energy \\
SU--XOM     & 25.86 & Energy \\
AMAT--LRCX  & 25.13 & Technology \\
\bottomrule
\end{tabular}
\label{tab:permutation}
\end{table}

Airline pairs (AAL, DAL, UAL) and energy pairs (HAL, SLB, KMI, WMB, CVX, XOM) dominate the most statistically significant results. The semiconductor pair AMAT--LRCX (from Phase~1) also achieves strong significance ($Z = 25.13$), confirming it generalizes across the larger universe.

\subsection{Multiple Testing Considerations}

With 3,539 pairs tested, the 32.4\% significance rate is well above the 5\% expected under the null hypothesis. Under random clustering, $\sim$177 pairs would exceed $Z > 1.645$ by chance; the observed 1,148 is 6.5$\times$ this, indicating genuine signal. However, no formal multiple-testing correction (e.g., Benjamini-Hochberg FDR) was applied. The top pairs with $Z > 20$ are far beyond any reasonable false-discovery threshold; marginal pairs near $Z = 1.65$ should be treated with more caution.

Additionally, the permutation test uses 30 permutations per timestamp, which yields a coarse null distribution estimate. Standard practice recommends $\geq 100$ permutations; the lower count was chosen as a computational concession given the cost of re-running the full OPTICS pipeline per permutation.

% ============================================================================
% 8. DISCUSSION
% ============================================================================
\section{Discussion}

\subsection{Phase 1 vs Phase 2 Comparison}

\begin{table}[H]
\centering
\caption{Comparison across phases.}
\begin{tabular}{@{}lccc@{}}
\toprule
\textbf{Metric} & \textbf{Phase 1} & \textbf{Phase 2 (Baseline)} & \textbf{Phase 2 (Kalman)} \\
\midrule
Universe        & 40 tickers, 1 sector & 142 tickers, 5 sectors & Same \\
Pairs tested    & $\sim$33             & 3,643                  & 50 (top by freq) \\
Profitable OOS  & 54\%                 & 54\%                   & 64\% \\
Top Sharpe      & 3.55                 & 2.56                   & 18.47 \\
Validation      & 3-test binary        & 5-test scored          & Same \\
Transaction costs & None               & None (baseline)        & 10bps \\
\bottomrule
\end{tabular}
\label{tab:phase-comparison}
\end{table}

Profitability is comparable ($\sim$54\%) for the baseline across phases, suggesting the framework generalizes. The Kalman strategy pushes this to 64\%. The 5-test scored framework dramatically increases the number of tradeable pairs (59.0\% vs $\sim$0\% with cointegration).

\subsection{Framework Evolution}

The progression from Phase~1 to Phase~2 addressed every limitation identified in Phase~1's Future Work:

\begin{enumerate}
    \item \textbf{Walk-forward validation:} Implemented with 5-split rolling windows.
    \item \textbf{Adaptive hedge ratios:} Kalman terminal beta estimation replaces static OLS.
    \item \textbf{Expanded universe:} 142 tickers across 5 sectors via systematic screening.
    \item \textbf{Transaction costs:} 10bps round-trip incorporated.
\end{enumerate}

\subsection{Strengths}

\begin{enumerate}
    \item \textbf{Methodological rigour:} Proper cal/OOS split, noise-adjusted frequencies, permutation tests, walk-forward validation, and the Kalman spread fix eliminate common sources of bias.
    \item \textbf{Multi-phase validation:} Testing on semiconductors first, then generalizing across sectors, provides cross-validation of the entire approach.
    \item \textbf{Intellectually honest:} We report that cointegration fails (0\%), that consensus does not predict performance, and that trade count is the binding constraint.
    \item \textbf{Genuine discoveries:} Semiconductor subsectors, crypto mining cluster, and RTX--SHEL emerge from raw price data alone.
    \item \textbf{Statistical foundation:} 1,148 permutation-significant pairs confirm that clustering structure is non-random.
\end{enumerate}

\subsection{Limitations}

\begin{enumerate}
    \item \textbf{Short OOS period ($\sim$68 days):} Most pairs generate 2--3 trades, below the 5-trade minimum for meaningful Sharpe ratios.
    \item \textbf{$\sim$12 months of hourly data:} Limited historical depth reduces ADF power and limits the optimizer's training set.
    \item \textbf{Simplified transaction costs:} Flat 10bps does not capture variation in spreads, market impact, or short-borrowing costs.
    \item \textbf{No portfolio effects:} Each pair evaluated independently. Correlated pairs (many energy pairs share tickers) require portfolio-level risk management.
    \item \textbf{Cross-sector pairs are weaker:} 55.6\% tradeable rate vs 64.1\% intra-sector. Most cross-sector relationships have lower co-clustering frequency, suggesting they are more transient by nature.
    \item \textbf{Fixed Kalman transition covariance (0.01):} Not optimized per pair.
\end{enumerate}

% ============================================================================
% 9. FUTURE WORK
% ============================================================================
\section{Future Work}

Items marked with $\checkmark$ have been completed in Phase~2.

\begin{enumerate}
    \item[$\checkmark$] \textbf{Walk-forward validation:} 5-split rolling cal/OOS windows.
    \item[$\checkmark$] \textbf{Adaptive hedge ratios:} Kalman terminal beta estimation.
    \item[$\checkmark$] \textbf{Expanded universe:} 142 tickers, 5 sectors via 3-layer screening.
    \item[$\checkmark$] \textbf{Transaction cost modelling:} 10bps round-trip incorporated.
    \item \textbf{Transient event trading (next priority):} Trade cluster formation/dissolution events directly---enter when a pair enters a cluster, exit on dissolution. This is the core application of the clustering methodology. The formation/dissolution detection infrastructure already exists in the codebase; the remaining work is wiring it into the backtest as trade signals and evaluating performance on short-lived relationships.
    \item \textbf{Real-time pipeline:} Live clustering $\to$ formation detection $\to$ signal generation.
    \item \textbf{Portfolio optimization:} Correlated pairs risk management and position sizing.
    \item \textbf{Longer historical data:} Multi-year hourly data for more robust parameter estimation and longer OOS windows.
    \item \textbf{Regime-aware position sizing:} Adjust exposure based on market volatility regime.
\end{enumerate}

% ============================================================================
% 10. CONCLUSION
% ============================================================================
\section{Conclusion}

We demonstrate that unsupervised clustering on hourly price-derived features is an effective pair discovery mechanism that generalizes beyond a single sector. In Phase~1, three algorithms converge on 8 consensus semiconductor pairs reflecting genuine subsector structure, with 54\% profitable out of sample. In Phase~2, the framework extends to 142 tickers across 5 sectors, yielding 3,643 scored pairs under a new 5-test validation framework that correctly accommodates transient (non-cointegrated) correlations.

Key contributions include:

\begin{enumerate}
    \item \textbf{5-test scored framework:} Replaces cointegration (0\% pass rate) with a system that classifies 59.0\% of pairs as tradeable while maintaining statistical rigour.
    \item \textbf{Kalman hedge with terminal beta:} Documents and fixes the artificial mean-reversion pitfall of adaptive Kalman spreads; the fixed-beta approach achieves 64\% profitability with transaction costs.
    \item \textbf{Statistical significance:} Permutation testing confirms 1,148 of 3,539 pairs have non-random clustering structure.
    \item \textbf{Cross-sector discovery:} The crypto mining cluster (APLD/HUT/CIFR/IREN) and RTX--SHEL emerge from price data alone, demonstrating the algorithm's ability to detect economically meaningful relationships across sector boundaries.
    \item \textbf{Honest reporting:} Consensus does not predict OOS performance, trade count is the binding constraint, and cross-sector pairs are weaker than intra-sector---all documented as important negative findings.
\end{enumerate}

The framework reliably identifies genuine relationships and produces partially profitable trading signals. Crucially, the current backtesting approach---trading persistent pairs with a static z-score strategy over the full OOS window---serves as a \textit{validation exercise} for the clustering methodology, not as the intended end application. The pairs that repeatedly co-cluster (airlines, energy majors, equipment stocks) confirm that the discovery engine finds real structure, but the unique edge of this research lies elsewhere.

The natural next step is \textbf{transient event trading}: entering positions when a pair forms a cluster and exiting when it dissolves, rather than trading continuously. Transient correlations can form and dissolve within hours or days. The formation/dissolution detection infrastructure already exists in the codebase. An event-driven strategy would trade \textit{with} the clustering signal---exploiting brief, surprising formations that traditional pairs traders are not monitoring---rather than applying a static strategy to pairs that happen to cluster frequently. The statistical foundation established here (permutation-validated clustering, 5-test scored validation, Kalman hedge ratios) provides the methodological groundwork for that future application.

% ============================================================================
% REFERENCES
% ============================================================================
\begin{thebibliography}{10}

\bibitem{engle1987}
Engle, R.F. and Granger, C.W.J. (1987).
\textit{Co-integration and error correction: Representation, estimation, and testing.}
Econometrica, 55(2), 251--276.

\bibitem{gatev2006}
Gatev, E., Goetzmann, W.N. and Rouwenhorst, K.G. (2006).
\textit{Pairs trading: Performance of a relative-value arbitrage rule.}
The Review of Financial Studies, 19(3), 797--827.

\bibitem{ankerst1999}
Ankerst, M., Breunig, M.M., Kriegel, H.-P. and Sander, J. (1999).
\textit{OPTICS: Ordering points to identify the clustering structure.}
ACM SIGMOD Record, 28(2), 49--60.

\bibitem{lomackinlay1988}
Lo, A.W. and MacKinlay, A.C. (1988).
\textit{Stock market prices do not follow random walks: Evidence from a simple specification test.}
The Review of Financial Studies, 1(1), 41--66.

\bibitem{hurst1951}
Hurst, H.E. (1951).
\textit{Long-term storage capacity of reservoirs.}
Transactions of the American Society of Civil Engineers, 116, 770--799.

\bibitem{kalman1960}
Kalman, R.E. (1960).
\textit{A new approach to linear filtering and prediction problems.}
Journal of Basic Engineering, 82(1), 35--45.

\bibitem{chan2013}
Chan, E.P. (2013).
\textit{Algorithmic Trading: Winning Strategies and Their Rationale.}
John Wiley \& Sons.

\bibitem{ester1996}
Ester, M., Kriegel, H.-P., Sander, J. and Xu, X. (1996).
\textit{A density-based algorithm for discovering clusters in large spatial databases with noise.}
KDD, 96(34), 226--231.

\bibitem{vidyamurthy2004}
Vidyamurthy, G. (2004).
\textit{Pairs Trading: Quantitative Methods and Analysis.}
John Wiley \& Sons.

\bibitem{dickey1979}
Dickey, D.A. and Fuller, W.A. (1979).
\textit{Distribution of the estimators for autoregressive time series with a unit root.}
Journal of the American Statistical Association, 74(366a), 427--431.

\end{thebibliography}

\end{document}
